\section{Conclusion}
\label{chap:conclusion}

\subsection{Synthèse}
Ce mémoire a présenté la conception, l’implémentation et l’évaluation de \nsit, un moteur de \acrlong{vn} web fondé sur un format de script déclaratif en \gls{json}. L’idée de départ était simple: vérifier si une architecture de ce type pouvait réellement permettre à des auteurs non-développeurs de créer et de publier des récits interactifs sans passer par une installation locale ou une phase de compilation.

Au fil du travail, j’ai défini un \gls{schema} de script, construit un \glsit{runtime} web en \gls{ts}/\gls{reactjs}, optimisé les composants les plus coûteux, notamment \texttt{Typewriter}, puis ajouté une validation d’exécution via \gls{zod}. Trois itérations successives, centrées sur les bases, l’optimisation et la robustesse, ont permis d’atteindre les objectifs classés \textit{Must} et \textit{Should}: affichage progressif du texte, navigation entre scènes, gestion des choix, validation du format et gestion basique des saisies. Certaines fonctions \textit{Can}, comme les \glsplit{sprite}, les transitions ou la sauvegarde déléguée à un \gls{webworker}, ont également été amorcées.


\subsection{Réponse à la problématique et bilan critique}
La question de recherche portait sur la capacité d’une architecture de moteur de \gls{vn} web, fondée sur un format déclaratif, à permettre à des auteurs non-développeurs de créer du contenu interactif. Sur le plan technique, la réponse est positive: j’ai montré qu’il était possible de décrire un récit interactif dans un fichier déclaratif et de l’exécuter directement dans un navigateur, avec des mécanismes de validation et des messages d’erreur exploitables.

En revanche, l’évaluation met aussi en évidence une limite plus nette du point de vue de l’auteur. Le format \gls{json}, même s’il est adapté au moteur, reste sensible aux erreurs de syntaxe et peu confortable pour une personne qui ne vient pas du développement. Sans éditeur guidé, validation en temps réel ou modèles prêts à l’emploi, l’approche déclarative seule ne suffit pas à rendre la création totalement autonome. Par ailleurs, certains choix techniques, comme la précomputation des \glslink{snaphtml}{snapshots} ou la délégation de traitements à un \gls{webworker}, améliorent le comportement du moteur, mais introduisent aussi des compromis qu’il faut accepter selon le contexte d’usage.

En résumé, l’architecture est viable et propre sur le plan technique, mais elle doit encore être entourée d’outils pensés pour l’écriture afin de répondre pleinement à l’objectif initial.


\subsection{Pistes de développement}

\subsubsection{Pistes de développement du Media Project}
Pour rendre la solution réellement utilisable par des auteurs non-développeurs, plusieurs évolutions me paraissent prioritaires :

\begin{itemize}
    \item Développer un éditeur guidé avec validation en temps réel (prévisualisation, messages d'erreur explicites, modèles), afin de masquer la complexité syntaxique du \gls{json}.
    \item Ajouter des outils de conversion et d'import (par ex. importer depuis un format plus lisible ou depuis un langage narratif type Ink) pour réduire la courbe d'apprentissage.
    \item Étendre le moteur avec des embranchements conditionnels, des variables de scénario plus riches et une API d'extensions pour permettre des fonctionnalités avancées sans modifier le cœur du \glsit{runtime}.
    \item Automatiser la chaîne de publication (intégration continue → déploiement statique) et fournir des \textit{workflows} \textit{one-click}\footnote{\textbf{Workflow} « one-click »: enchaînement automatisé d'étapes techniques (ici, de compilation et de déploiement) déclenché par une seule action de l'utilisateur.} pour la mise en ligne des histoires.
    \item Améliorer l'accessibilité et la documentation utilisateur (guides pas-à-pas, exemples, galerie de templates).
\end{itemize}

\clearpage


\subsubsection{Pistes de recherche}
Plusieurs axes de recherche découlent de ce travail :

\begin{itemize}
    \item L'ajout d'un éditeur graphique ou assisté, visuel ou nodal, superposé au format déclaratif, qui masquerait la syntaxe \gls{json} tout en conservant les avantages du texte brut, comme le suivi de version ou la portabilité.
    \item La conception d'un langage narratif intermédiaire, plus proche du langage naturel, qui serait ensuite compilé vers le \gls{schema} \gls{json} du moteur, à la manière de langages dédiés comme Ink ou Twine.
    \item L'intégration d'assistants d'écriture fondés sur des modèles de langage, capables de générer ou de restructurer des fragments narratifs directement dans le format déclaratif.
    \item L'extension du modèle déclaratif vers une logique narrative plus riche, avec des variables d'état persistantes ou des embranchements conditionnels, et son effet sur la lisibilité du format pour un auteur non-développeur.
    \item La standardisation ou l'interopérabilité d'un format déclaratif de \gls{vn}, de façon à ce qu'un même script puisse être exécuté par plusieurs moteurs ou plateformes.
    \item La prise en compte de l'accessibilité, notamment pour les lecteurs d'écran et l'internationalisation, directement au niveau du schéma déclaratif plutôt que comme une couche ajoutée après coup.
\end{itemize}

Ces pistes ouvrent des perspectives complémentaires à ce travail, en interrogeant la manière dont le format déclaratif pourrait évoluer pour concilier simplicité d'accès et richesse narrative.